% Transcription — fonds Alexandre Grothendieck, shelfmark 158, batch 3 (pages 41-60).
% Established from the facsimile archives/batches/158/batch-03.pdf (a mirror of
% https://grothendieck.umontpellier.fr/158.pdf), per the skill
% transcribe-grothendieck. No cover sheet inside the batch file: batch PDF page
% N = archivists' page 40+N, checked against the pencilled numbers, which sit
% bottom-left on most leaves and bottom-right on page 41.
% Pages 43, 54 and 59 are blank and carry nothing but the archivists' number.
% Page 42 is an unrelated administrative verso, scanned upside down — absent.
% Pass: Opus 5 (claude-opus-5), 2026-09-20 — first pass, unchecked against the pages by a human.
\documentclass[11pt,a4paper]{article}
% Le préambule commun à toutes les transcriptions du fonds Grothendieck.
%
% Il tient en une idée : **l'incertitude se compose, elle ne se résout pas.**
% Une transcription qui lisse un mot douteux a détruit la seule chose pour
% laquelle on la faisait — un lecteur doit pouvoir savoir, à chaque endroit,
% ce qui était sur le feuillet et ce qui a été ajouté. Les six macros qui
% suivent sont donc l'appareil critique, et non de la décoration.
%
% Les mêmes commandes sont reconnues par `scripts/render.mjs`, qui produit la
% vue de lecture du site. Une macro ajoutée ici doit l'être aussi là-bas,
% faute de quoi le rendu échouera bruyamment — ce qui est voulu : un rendu
% silencieusement faux est pire qu'un rendu absent.

\ProvidesPackage{grothendieck}[2026/08/08 Transcriptions du fonds Grothendieck]

\usepackage{amsmath,amssymb,amsthm}
\usepackage{mathrsfs}
\usepackage{stmaryrd}
% Les diagrammes commutatifs sont partout dans ce fonds : c'est la langue
% même de la géométrie algébrique de Grothendieck, et une transcription qui
% les rendrait en prose ne transcrirait rien.
\usepackage{tikz-cd}
% Français par défaut — c'est la langue des feuillets — mais l'anglais est
% chargé aussi : la traduction et le résumé sont des documents anglais, et la
% typographie française (espaces avant les deux-points) y serait une faute.
% Ces documents basculent par \selectlanguage{english} en tête de corps.
\usepackage[english,french]{babel}
\usepackage{fontspec}
\usepackage{geometry}
\geometry{a4paper,margin=25mm,marginparwidth=28mm}
\usepackage{marginnote}
\usepackage{xcolor}
% ulem, not soul. `soul`'s \st and \ul are built on a character-by-character
% scan that XeLaTeX's font machinery breaks: the document fails with an
% undefined control sequence rather than degrading. ulem does the same two jobs
% and is Unicode-engine safe. `normalem` stops it hijacking \emph, which the
% transcriptions use for Grothendieck's own emphasis.
\usepackage[normalem]{ulem}
\usepackage{hyperref}
\hypersetup{colorlinks=true,linkcolor=black,urlcolor=black,pdfborder={0 0 0}}

% --- Métadonnées du lot -----------------------------------------------------
% Reprises telles quelles par le script de rendu, qui les affiche en tête de
% la vue de lecture. Elles fixent ce que le fichier prétend couvrir : sans
% elles, une transcription flotte sans point d'attache dans un fonds de seize
% mille feuillets.

\newcommand{\folder}[1]{\gdef\grothFolder{#1}}
\newcommand{\batch}[1]{\gdef\grothBatch{#1}}
\newcommand{\leaves}[2]{\gdef\grothFirst{#1}\gdef\grothLast{#2}}
\newcommand{\foldertitle}[1]{\gdef\grothFoldertitle{#1}}
\gdef\grothFolder{?}\gdef\grothBatch{?}\gdef\grothFirst{?}\gdef\grothLast{?}\gdef\grothFoldertitle{}

% --- Le résumé --------------------------------------------------------------
%
% Il ouvre la lecture modernisée, sous le titre « Résumé », et il est composé
% en petit. Ce n'est pas un ornement : c'est le seul passage du document qui
% ne soit pas des mathématiques, et un lecteur doit pouvoir le reconnaître
% avant de l'avoir lu — puis le sauter s'il connaît déjà le sujet, ou s'y
% arrêter s'il ne le connaît pas. Le filet qui le suit marque la reprise du
% corps.

\newenvironment{resume}
  {\small\setlength{\parskip}{0.45em}}
  {\par\medskip\hrule\medskip}

% --- Mots-clés ---------------------------------------------------------------
%
% En anglais, en clôture du résumé de la lecture modernisée : le vocabulaire
% moderne sous lequel chercher ce que les feuillets construisent. Le manifeste
% du site (`npm run manifest`) les extrait pour étiqueter la cote entière —
% cette ligne est la source unique des tags, il n'y a pas de fichier à côté.
\newcommand{\keywords}[1]{\par\smallskip\noindent
  {\footnotesize\textsc{Keywords} --- \textit{#1}}\par}

% --- L'appareil critique ----------------------------------------------------

% Le numéro de feuillet, en marge. C'est l'ancre qui permet de régler un
% désaccord sur une formule en regardant *un* feuillet plutôt qu'une chemise
% de six cents. Le site s'en sert aussi pour tourner les pages du fac-similé
% quand on fait défiler la transcription.
\newcommand{\leaf}[1]{%
  \marginnote{\footnotesize\color{gray!70}\textsf{#1}}%
  \label{leaf:#1}\ignorespaces}

% Illisible : on ne devine pas. Un mot inventé qui se lit comme les autres est
% le pire résultat possible.
\newcommand{\ill}{\textcolor{red!60!black}{[\,\ldots\,]}}

% Lecture douteuse : le mot est donné, et signalé comme incertain.
\newcommand{\uncertain}[1]{\uline{#1}}

% Ajout de l'éditeur — une lettre manquante, un mot sous-entendu.
\newcommand{\add}[1]{\textcolor{blue!50!black}{[#1]}}

% Biffé par Grothendieck lui-même. Ce qu'il a rayé fait partie du document :
% une rature dit souvent où la pensée a changé de direction.
\newcommand{\struck}[1]{\sout{#1}}

% Note du transcripteur, sur l'état matériel du feuillet.
\newcommand{\note}[1]{\par\smallskip{\small\color{gray!60!black}\textit{[#1]}}\par\smallskip}

% Note marginale de Grothendieck, distincte du corps du texte.
\newcommand{\marginal}[1]{\par\smallskip\noindent\hspace{1em}%
  {\small\color{gray!60!black}#1}\par\smallskip}

% --- Titre ------------------------------------------------------------------

\AtBeginDocument{%
  \begin{center}
    {\large\bfseries Fonds Alexandre Grothendieck --- cote n\textsuperscript{o}~\grothFolder}\\[2pt]
    {\itshape\grothFoldertitle}\\[4pt]
    {\small feuillets \grothFirst--\grothLast\ (lot \grothBatch)}\\[2pt]
    {\footnotesize\color{gray!60!black}Transcription établie d'après le fac-similé numérisé
     par l'Université de Montpellier.\\
     Les lectures incertaines sont soulignées, les illisibles notées [\,\ldots\,],
     les ajouts entre crochets.}
  \end{center}
  \vspace{1em}}

\endinput


\folder{158}
\batch{3}
\pages{41}{60}
\dating{[vers 1990-1991]}
\watermark{Édition de démonstration}
\foldertitle{[Autour des Dérivateurs] : notes manuscrites (s.d.)}

\begin{document}

\section*{Monoïdes finis, catégories pseudo-cofiltrantes et 1-connexité}

\note{les intertitres de cette transcription sont de l'éditeur : aucun des
feuillets du lot n'en porte. Le lot réunit trois suites paginées de sa main,
indépendantes les unes des autres : la fin d'une première suite (feuillet 41,
numéroté 7), une suite complète numérotée 1 à 8 (feuillets 44 à 53), et une
suite numérotée 1 à 3 (feuillets 55 à 57), suivie de feuillets d'essai}

\note{conventions de notation retenues pour tout le lot : $X^\wedge$ pour son
$X$ surmonté d'un accent circonflexe (la catégorie des préfaisceaux sur $X$),
qu'il trace tantôt en exposant, tantôt superposé à la lettre ; $\mathbb{N}$ et
$\mathbb{Z}$ pour ses $N$ et $Z$ barrés ; $X_{/\!/a}$ pour son $X$ à double
barre en indice}

\page{41}
\note{la page porte « 7 » de sa main en tête ; elle commence au milieu d'une
phrase et fait suite, pour le sens, à ce qui s'interrompt au bas du lot
précédent. Sa pagination ne s'y raccorde pourtant pas : les feuillets 21 à 39
portent de sa main 3 à 18, de sorte qu'un « 7 » ne peut clore cette suite-là.
On laisse le chiffre tel qu'il est lu et l'on signale le désaccord plutôt que
de le résoudre}

$u \in M$ et pour $p, q \in \mathbb{N}$, $\varphi^p u = \varphi^q u \Rightarrow
p = q$, on aura \struck{que $\varphi^p$} (on \uncertain{remarquera} \ill{} on
$u = \varphi$ puisque \struck{$\exists$} $u'$ avec $uu' = \varphi$)
$\varphi^p = \varphi^q \Rightarrow p = q$.

Dans ce cas, par des principes généraux, \struck{le quotient} l'objet $e$ de
$X^\wedge$ étant 0-connexe, l'objet quotient $F = e/\mathbb{N} = e/\varphi$ est
0-connexe et son $\pi_1$ admet un quotient isomorphe \uncertain{à}
$\mathbb{Z}$, donc $X^\wedge/F$ n'est pas 1-connexe, et
$H^k(F, \mathbb{Z}) \neq 0$ pour \uncertain{un entier} $k \neq 0$.

b) Il existe $p \neq q$ tel que $\varphi^p = \varphi^q$. Alors
\struck{l'opération} l'image de $M$ dans l'ens. $\mathrm{Ep}(E)$ des
épimorphismes de $E$ $=$ ens. des éléments minimaux de $E$ est un groupe
$\Gamma$, et si $X_0$ désigne l'ouvert de $X$ formé des plus petits objets de
$X$, alors \struck{\ill{}}
\[
  \pi_1(X_0, e) \;\simeq\; \pi_1(M) \twoheadrightarrow \Gamma .
\]
Donc si $X_0$ est 1-connexe, alors $\Gamma = \{1\}$, donc \ill{}
$u\varphi = \varphi$ $\forall u \in M$, et $M$ \struck{est} \ill{} donc ainsi
$X_0$, \struck{\ill{}} est cofiltrant.

\note{tel qu'il est écrit, « des épimorphismes de $E$ $=$ ens. des éléments
minimaux de $E$ » définit $E$ par lui-même ; l'argument appelle un $E_0$,
l'ensemble des éléments minimaux de $E$, qui reparaît sous ce nom au
feuillet 45}

\textbf{Conclusion} Si tous les objets de $X^\wedge$ \ill{} est 1-connexe,
alors $X$ est cofiltrant.

\marginal{\ill{} $\mathrm{End}(E)$ \ill{} $\varphi = \varphi$ $X_0$
\struck{\ill{}} $\pi_1(\ill{})$ \ill{}}
\note{note marginale portée en diagonale dans l'angle inférieur gauche, en
grande partie illisible}

\subsection*{Proposition 1 : un monoïde fini pseudo-cofiltrant à groupe enveloppant trivial}

\page{44}
\note{la page porte « 1 » de sa main : début de la suite qui court jusqu'au
feuillet 53}

\textbf{Proposition 1} Soit $M$ un monoïde fini, tel que

a) le groupe enveloppant de $M$ soit $\simeq \{1\}$, et

b) $M$ est pseudo-cofiltrant. \struck{pseudo-cofiltrant.}

Alors il existe un \struck{\ill{}} élément $p$ de $M$ tel que l'on ait, pour
tout $u \in M$
\[
  (*) \qquad u p \;\struck{= up}\; = p , \qquad
  \text{(en particulier } p^2 = p\text{)},
\]
\struck{en particulier} donc $M$ est \struck{filtrant} cofiltrant.

\textbf{Dém.} \struck{L'unicité de $p$ est claire, car si $p' \in M$
\ill{} satisfait aux mêmes propriétés que $p$, on a (faisant $u = p'$ dans
$(*)$) $p'p = pp' = p$ et symétriquement \ill{} $pp' = p'pp' = p'$, donc
$p = p'$.}
\note{tout ce paragraphe de démonstration est barré de deux longues
diagonales ; il se lit néanmoins d'un bout à l'autre}

\struck{On peut supposer} $M$ pseudo-cofiltrant \struck{\ill{}} implique que
pour un \struck{ensemble} nb fini d'\struck{éléments} objets $u_i$ de $M$, il
existe un \struck{\ill{}} objet $p$ de $M$ tel que \struck{pour tout} pour tout
$i$, $p$ se factorise en
\[
  p \;=\; u_i v_i .
\]
Comme $M$ est fini il existe donc un $p$ tel que pour tout $u \in M$, on ait
une factorisation
\[
  (**) \qquad p \;=\; u\, v(u) ,
\]
où $v(u)$ est un objet de $M$. \struck{Cette propriété} L'existence d'un tel
$p$ implique visiblement que $M$ est cofiltrant, donc équivaut (pour $M$ fini)

\page{45}

à l'hypothèse b).

OPS $M$ réalisé comme monoïde d'endomorphismes d'un ensemble \struck{fini} $E$
(p.~ex. par la représentation régulière gauche de $M$).
\note{« OPS » = on peut supposer ; l'abréviation revient dans le lot}

La relation $(**)$ implique
\[
  \mathrm{Im}\, p \;\subset\; \mathrm{Im}\, u \qquad \forall u \in M,
\]

\marginal{$E_0 =$}
\note{la mention ci-dessus est portée en interligne, une accolade la rattachant
au $\mathrm{Im}\, p$ de la phrase suivante}
donc $\mathrm{Im}\, p$ est le plus petit des $\mathrm{Im}\, u$, pour
$u \in M$. Soit $u \in M$, on a
\[
  u(E_0) \;=\; u\, p(E) \;=\; \mathrm{Im}(up) \;\supset\; E_0 ,
\]
et comme $E_0$ est fini, on en conclut que $u(E_0) = E_0$ et que $u|E_0$ est un
\uncertain{automorphisme} de l'ens. $E_0$. Soit $u_{E_0}$ cet
\uncertain{automorphisme}. L'application
\[
  u \longmapsto u_{E_0} \;:\; M \longrightarrow \mathrm{Aut}(E_0)
\]
est un hom. de $M$ dans un groupe. L'hyp. a) sur $M$ implique donc qu'on a
$u_{E_0} = \mathrm{id}_{E_0}$ $\forall u \in M$. On en conclut aussi que
\[
  E_0 \;=\; \text{ens. des pts fixes de } M \text{ opérant dans } E,
\]
car si $x \in E \setminus E_0$, comme $p(x) \in E_0$, on a $p(x) \neq x$

\page{46}

donc $x$ n'est pas pt fixe pour $M$. \struck{On en conclut, d'où}
$\forall u \in M$ la relation $u_{E_0} = \mathrm{id}_{E_0}$ s'écrit \ill{}
\[
  u\, p(x) \;=\; p(x) \quad \forall x \in E, \qquad \text{i.e.}
\]
\[
  up \;=\; p \qquad \forall u \in M,
\]
ce qui prouve déjà que $M$ est cofiltrant (car $u, v \in M$, on a $up = vp$,
les deux \struck{membres} étant égaux à $p$), cqfd.

J'ai envie de proposer ceci :

\textbf{Proposition 2 (??)} \add{S}oit $X$ une catégorie 1-connexe
\emph{finie}, ayant un \struck{objet} plus grand objet $a$), et
pseudo-cofiltrante. Alors $X$ est cofiltrante.

On retrouve $\uncertain{(\ddagger)}$ le \ill{}~1 en prenant $X$ réduite à un
seul objet. Je vais essayer de calquer la dém., au \ill{}~2. Je note
\struck{qu'il est cat. \emph{finie}}, \ill{} dans la catégorie $X/a$, (il y a
un plus petit objet,

\note{le mot qui précède « 1 », puis « 2 », n'est pas lisible ; le sens appelle
« la proposition », énoncée plus haut sous ce numéro}

\page{47}
\note{la page porte « 4 » de sa main}

soit $(e_a, p_a)$, ou \ill{}
\[
  p \;:\; e \longrightarrow a .
\]
Si $b$ est dans $X$, on voit que la \uncertain{dernière, dans} catégorie
$X/b$ il y a un plus petit \struck{\ill{}} objet, et on voit
\struck{facilement} \struck{\ill{}} que \struck{on} peut prendre à source égale
à $e$, donc de la forme
\[
  e \xrightarrow{\;p_b\;} b .
\]
Mais \struck{\ill{}} il n'y a pas unicité de $p_b$, et pour \ill{} une
\ill{}, à faire ce point, on voudrait mettre en relation le choix de $p_b$,
avec celui (déjà fait) de $p_a = p$. Comme $a$ est un plus grand objet de $X$,
il existe
\[
  u \;:\; b \longrightarrow a
\]
et on voit facilement \ill{} que pour $u$ fixé, on peut choisir $p_b$ de la
façon \ill{} que

\begin{tikzcd}
& b \arrow[dr, "u"] & \\
e \arrow[ur, "p_b"] \arrow[rr, "p"'] & & a
\end{tikzcd}

commutatif.

\page{48}
\note{feuillet d'essai : diagrammes repris deux fois et formules isolées, sans
prose suivie. On transcrit ce qui se lit}

$p \in \mathrm{Hom}(e,a)$

\begin{tikzcd}
& b \arrow[dr, "u", bend left=15] \arrow[dr, "v"', bend right=15] & \\
e_a \arrow[ur, bend left=15] \arrow[ur, bend right=15]
    \arrow[rr, "p_a"'] & & a
\end{tikzcd}

\note{les deux flèches de $e_a$ vers $b$ sont notées $p_{b,u}$ et $p_{b,v}$ sur
le feuillet}

$\varphi \in \mathrm{Hom}(e,e)$

\begin{tikzcd}
e \arrow[r, "p"] \arrow[d, bend right=25] & a \\
e \arrow[u, bend right=25] \arrow[ur] &
\end{tikzcd}

\[
  uq = p, \qquad \struck{\ill{}} \text{ avec } q \text{ comme } p,
  \qquad \text{donc } uq = q
\]
\[
  \underbrace{u\,q}_{q} \;=\; \underbrace{v\,q'}_{q'} \;=\; p
\]

\begin{tikzcd}
& b \arrow[dr, "u"] & \\
e \arrow[ur, "p_u"] \arrow[rr] & & q
\end{tikzcd}

\begin{tikzcd}
& q & \\
e \arrow[ur] \arrow[dr, "p_v"'] & & \\
& b \arrow[uu, "v"'] &
\end{tikzcd}

\struck{$\mathrm{Hom}(x,a)$}, $\mathrm{Hom}(m,a)$,
$x \mapsto \mathrm{Hom}(e,x)$

\[
  up = q, \qquad p = v\,\overline{v}, \qquad p = q\,q', \qquad
  q = v\,\struck{\ill{}}
\]

\subsection*{Corollaire : $X$ finie, pseudo-cofiltrante, $U$ 1-connexe}

\page{49}
\note{la page porte de sa main un chiffre corrigé, illisible sous la
surcharge}

\textbf{Corollaire} Soit $X$ une catégorie \emph{finie},
(i.e. $X^\wedge$ est loc. \struck{\uncertain{irréductible}}),
pseudo-cofiltr. \struck{sép.} \struck{Donc} $\mathrm{Ord}\, X$ est un ensemble
ordonné cofiltrant, admettant donc un plus petit élément. Soit $U$ l'ouvert de
$X$ formé des objets de $X$ qui sont des plus petits objets.
\struck{On suppose} \ill{} $U$ 1-connexe, \struck{alors} \struck{$X$} est
cofiltrante. (Donc si tout objet \struck{\ill{} de} $X$ est 1-connexe, alors
$X$ est cofiltr.)

En effet, il est immédiat \struck{que} \add{$U$ est} (et $X$ aussi)
pseudo-cofiltrante, un ouvert cocofinal dans $X$, que $X$ est cofiltrante dès
que $U$ l'est. D'autre part si $e \in U$, alors \add{$U_0$} l'inclusion de la
sous-catégorie pleine \add{(}de $X$ réduite\add{)} à $e$ dans $U$ est une
équivalence, ce qui implique que $U_0$ est 1-connexe, $U$ l'est
\uncertain{aussi}, et que $U$ est cofiltrant ssi $U_0$ l'est. Donc [\,il est
immédiat \uncertain{aussi}, \struck{$X$} étant pseudo-cofiltrant, que l'ouvert
cocofinal $U$ l'est \uncertain{également}, donc $U_0$ (équivalent : $U$)
\uncertain{également}.\,]

Soit donc $M = \mathrm{End}_X(e)$. On peut appliquer à $M$ la proposition 1,
qui implique que $M$ est cofiltrant, ainsi $U_0 \simeq B_M$ est cofiltrante,
cqfd

\marginal{Faux ! revoir la \uncertain{démonstration} \ill{}}

\textbf{N.B.} Au lieu de l'hypothèse $X$ finie, il suffisait de supposer que
a) $X$ admet un plus petit objet $e$, et b) $\mathrm{End}_X(e)$ est fini.

\subsection*{Seconde mouture : la proposition 1 reprise par les $H_x$}

\page{50}
\note{la page porte « 5 » de sa main ; l'énoncé de la proposition 1 est repris
depuis le début, sous des hypothèses différentes}

\textbf{Proposition 1} Soit $X$ une catégorie telle que

a) $X$ pseudocofiltrante --- b) $X$ admet un plus petit objet $e$
(conséquence de \struck{a)} si $\mathrm{Ob}\, X$ fini)
\struck{objet minimal}

c) Pour tout $x \in \mathrm{Ob}\, X$, l'ens. des
sous-$\underbrace{\mathrm{Hom}(e,e)}_{M}$-ens. à droite de $\mathrm{Hom}(e,x)$
satisfait la condition minimale (ce qui est le cas si les $\mathrm{Hom}(e,x)$
sont finis). Alors

$1^{\circ})$ $\forall x \in \mathrm{Ob}\, X$, parmi les sous-$M$-ens. à gauche
de $\mathrm{Hom}(e,x) \overset{\text{déf}}{=} F(x)$ qui sont de la forme $fM$
(avec $f \in F(x)$) il y en a
un plus petit, \struck{il est formé de l'\ill{}} \ill{} des
$\lambda \in \mathrm{Hom}(e,x)$ \struck{de la forme} \struck{$\forall f \in
\mathrm{Hom}($} $\forall f \in F(x)$, $\exists u \in M$ \ill{}
$\lambda = fu$.

\marginal{soit $H_x = \bigcap_{f \in F(x)} fM$.}
\note{la définition ci-dessus est portée en interligne au-dessus de « il y en a
un plus petit »}

$2^{\circ})$ $H_x$ est donc un $M$-ens. à droite. De plus, pour tout
$f : x \to y$ \ill{}
\[
  H_y \;=\; f \cdot H_x .
\]
On a donc $\mathrm{card}\, H_y \leq \mathrm{card}\, H_x$ si $x \leq y$
($\overset{\text{déf}}{\Longleftrightarrow}$ si $\mathrm{Hom}(x,y) \neq
\emptyset$).

$3^{\circ})$ Supposons les $F(x)$ finis (cf. condition c) ci-dessus) et que $e$
soit \uncertain{aussi} objet maximal i.e. que $\mathrm{Ord}\, X$ soit réduit

\page{51}
\note{la page porte « 6 » de sa main, écrit sur un chiffre biffé}

à un élément. Alors pour tout $f$ \uncertain{flèche}
$f : x \to y$ dans $X$, \struck{$x \to f$}
\[
  \lambda \longmapsto f\lambda \;:\; H_x \xrightarrow{\ \sim\ } H_y
\]
est bijectif (car les $H$ \uncertain{forment} un syst. local sur $X$).

$4^{\circ})$ Supposons, \struck{\ill{}} en plus des conditions a) b) c) et
celle de $3^{\circ})$, que $X$ soit 1-connexe. Alors le système local des $H_x$
est \add{donc} constant, \struck{dessus} en particulier \add{i.e. dont}
l'opération de $M = \mathrm{Hom}(e,e)$ sur $H_x$ est \struck{\ill{}} et par
suite l'\uncertain{constante}

\struck{$\lambda \in H_x$, $u \in M \Longrightarrow \lambda u = \lambda$}

\struck{\uncertain{bijections}}
\[
  u \longmapsto fu \;:\; H_e \longrightarrow H_x
\]
définie par un $f \in F(x) : e \to x$ ne dépend pas du choix de $f$, donc
\[
  fu \;=\; gu \qquad \text{pour } u \in H_e, \quad
  f, g : e \rightrightarrows x .
\]

\page{52}
\note{la page porte « 7 » de sa main ; la proposition 2 est elle aussi reprise
depuis le début}

Comme corollaire, on a ceci

\textbf{Proposition 2} Soit $X$ une catégorie \emph{finie}
pseudo-cofiltrante, \struck{\ill{}} et soit $X_0$ la catégorie \add{est}
ouverte de $X$ formée des \struck{objets} \add{plus petits} \ill{} de $X$
(car les hyp. impliquent $X_0 \neq \emptyset$), \struck{donc} $X_0$ est
0-connexe). Supposons que $X_0$ soit 1-connexe. Alors $X$ est cofiltrante
(donc tout objet 0-connexe de $X^\wedge$ est $\infty$-connexe, et plus
généralement $\infty$-asphérique pour tout \struck{\ill{}}
\uncertain{localisateur fondamental} $W$ satisfaisant \ill{}
$W(\uncertain{\delta})$).

À fortiori :

\textbf{Corollaire} Soit $X$ une catégorie \emph{finie} telle que tout objet
0-connexe de $X^\wedge$ soit 1-connexe (p.~ex. telle que $X^0$ soit tot.
$W$-asphérique, si $W$ est un loc. fondamental tel que $W \subset W_1$). Alors
$X$ est cofiltrante, i.e. $X^0$ est filtrante (et tout objet de $X^\wedge$ est
$\infty$-connexe).

\note{l'ouvert des plus petits objets est noté $X_0$ dans l'énoncé de la
proposition et $X^0$, l'indice porté en exposant, dans celui du corollaire ;
la transcription garde les deux graphies}

\page{53}
\note{la page porte « 8 » de sa main et ne compte que trois lignes : fin de la
suite}

En effet, on voit \ill{} que l'hyp. implique que $X$ est pseudo-cofiltrante,
et on applique la proposition 2.

\subsection*{Deux exemples : les épimorphismes et les monomorphismes d'un ensemble infini}

\page{55}
\note{la page porte « 1 » de sa main : début d'une suite nouvelle}

I) Soit $E$ un ens. \uncertain{infini},
\[
  M \;=\; \mathrm{Ep}(E)
\]
le monoïde des endos surjectifs de $E$. On a les propriétés suivantes

$1^{\circ})$ $uf = vf \Longrightarrow u = v$
\marginal{donc $M$ non cofiltrant}

$2^{\circ})$ $M$ est pseudo-cofiltrant, et \struck{$=$} il existe des éléments
minimaux, i.e. des \struck{\ill{}} $w$ tels que $\forall u \in M$, $\exists
w_u \in M$ avec
\[
  \struck{p}\, w \;=\; u\, w_u
\]
(et on peut \struck{$=$} prendre $w_u$ \add{minimal}). (Mais bien sûr on n'a
pas \struck{\ill{}} $uw = w$ $\forall u \in M$ ; car $up = p$ impliquerait
$u = 1$ \ill{}). On n'a donc pas non plus $w^2 = w$, ni $u^2 = u$ pour aucun
$u \in M$ sauf $u = 1$.) [\,Les $w$ minimaux sont ceux pour lesquels
$\mathrm{card}\, w^{-1}(x) = \pi \overset{\text{déf}}{=} (\mathrm{card}\, E)$
pour tout $x \in E$.\,]

$3^{\circ})$ Si $w_0$ est \ill{} minimal, l'ens. des \ill{} \ill{} est égal à
l'ens. des $w = w_0\,\sigma$ \ill{} \ill{} dans $M$) \ill{} avec
$\sigma \in M^{*}$ (aux \ill{} \uncertain{cycles} \ill{})
[\,donc $M^{*} = \mathrm{Aut}(E)$\,].
\note{tout ce paragraphe est traversé par une insertion diagonale surchargée,
elle-même partiellement biffée ; le peu qui s'en lit est donné ci-dessus}

$4^{\circ})$ \ill{} dans $M$, \ill{} existe $w$ minimal et $f$ telle que
$f\,w = f$ (donc l'image $\overline{w}$ de $w$ dans le groupe enveloppant \add{$G$}
de $M$ est 1).

$5^{\circ})$ Soit $M^{*}_0 \subset M^{*}$ \struck{\ill{}} la partie formée des
$\sigma \in M^{*}$ \struck{telles} qu'il existe $f \in M$ avec
$f\,\sigma = f$.

\page{56}
\note{la page porte « 2 » de sa main}

Alors $M^{*}_0$ engendre le groupe $M^{*}$. Ceci implique que
\struck{\ill{}} l'image de $M^{*}$ dans $G$ est égale à $\{1\}$, donc, compte
tenu des $3^{\circ})$ et $4^{\circ})$, que pour tout objet minimal $w$ de $M$,
on a $\overline{w} = 1$. Compte tenu de $2^{\circ}$,
\[
  w \;=\; u\, w_u \qquad \text{avec } w, w_u \text{ minimaux}
\]
on en conclut $\overline{u} = 1$ $\forall u \in M$, i.e.

$6^{\circ})$ le \struck{\ill{}} groupe enveloppant $G$ de $M$ est $\{1\}$.

II) Soit toujours $E$ un ens. \uncertain{infini}, et
\[
  M \;=\; \mathrm{Mon}(E)
\]
le monoïde des endos injectifs de $E$. Alors $M$ satisfait l'ensemble des
\struck{\ill{}} propriétés duales \add{de I} de $1^{\circ}$ à $6^{\circ})$,
i.e. $M^0$ satisfait aux propriétés $1^{\circ}$ à $6^{\circ}$. Les objets
minimaux \struck{sont} de $M$, i.e. \struck{elles} \add{tels} que
$\forall u \in M$, on ait une factorisation
\[
  w \;=\; w\, u , \qquad
  \mathrm{card}\bigl(E - w(E)\bigr) \;=\; \pi
\]
sont ceux pour lesquels \ill{}

\page{57}
\note{la page porte « 3 » de sa main}

\struck{$X_{/\!/a} \cap X_{/\!/b}$}

Soit $X$ loc. irréductible [\,i.e. telle que les comp.
\struck{irréductibles} \add{connexes} soient irréductibles i.e. satisfassent à
\uncertain{$P_S * 1'$}\,]. Alors la relation
\[
  R(a,b) \;:\; X_{/\!/a} \cap X_{/\!/b} \neq \emptyset
\]
est une relation d'équivalence dans $\mathrm{Ob}\, X$, laquelle n'est autre que
celle \struck{\ill{}} définie par
\[
  \mathrm{Ob}\, X \longrightarrow \pi_0(X) .
\]
De plus, si $X_{/\!/a} \cap X_{/\!/b}$ \struck{$\neq \emptyset$} est non vide,
il est 0-connexe, \struck{\ill{} et donc $=$ \ill{}}
\struck{Si $X$ est 0-connexe, \ill{} pour} \add{il suffit} que $X$ soit
1-connexe, il suffit que les $X_{/\!/a}$ le soient.

\marginal{\ill{} $X_{/\!/a} \neq \emptyset$ \ill{} \uncertain{connexe} \ill{}}
\note{accolade et note portées en diagonale dans la marge de gauche, presque
entièrement illisibles}

\page{58}
\note{feuillet d'essai, sans prose suivie}

\[
  \varphi u = \varphi v \quad \text{avec } u, v, \varphi \ \uncertain{surj.}
  \quad \Longrightarrow \quad u = v
\]
\[
  \varphi u = \varphi \quad \Longrightarrow \quad u = \mathrm{id}\ ?
  \qquad \uncertain{\mathrm{Mon}}
\]

\begin{tikzcd}
G \arrow[r, "u"] & F \arrow[r, "\varphi"] & E
\end{tikzcd}

\[
  M \;=\; \mathrm{Ep}(E)
\]
\marginal{\struck{description} des \uncertain{épimorphismes} de l'ens. $E$ dans
lui-même.}

$M$ est ps. cofiltrant, et \ill{} un objet \ill{} $\forall u \in M$ minimal $w$
(tel que pour tout $u \in M$, $\exists u'$ avec $w u' = w$) --- il suffit de
prendre pour $w$ une surjection dont les fibres sont de cardinal \ill{}
$\mathrm{card}\, E$. $M$ n'est pas cofiltrant, car
\[
  u f = v f \quad \Longrightarrow \quad u = v .
\]
\ill{} \uncertain{groupe} associé : $M$ est-il trivial ? Je suppose que oui.
On \struck{doit} obtenir des relations non triviales dans le groupe
enveloppant $G$, appliquées aux triples $u, v, f$ \ill{} avec
\[
  f u = f v , \qquad u \neq v
\]
ce qui donne $\overline{u} = \overline{v}$ dans $G$.

\page{60}
\note{feuillet d'essai portant « (1) » cerclé de sa main : carquois tracés sans
étiquettes, formules isolées. On donne ce qui porte du texte ; les graphes non
étiquetés ne sont pas reproduits}

\begin{tikzcd}
F_s \arrow[r] & F_{s'}
\end{tikzcd}

\begin{tikzcd}
s & s' \arrow[l]
\end{tikzcd}

\begin{tikzcd}
X \arrow[r] & \mathrm{Groupoïdes}
\end{tikzcd}

2-catégorie des 2-\struck{groupes} \add{champs}

\struck{$X_0$}, \struck{$X_1$},
$\mathcal{F}_{0u} \xleftarrow{\ \sim\ } \mathcal{F}_{1t}$

\note{le diagramme qui suit est redessiné d'un croquis serré ; les deux
inclusions y portent la mention « fermé », et la flèche verticale de droite
descend d'un « groupoïde loc. constant »}

\begin{tikzcd}[column sep=small, row sep=small, nodes={font=\scriptsize}]
& Z \arrow[d] & \Psi \arrow[l, hook'] \arrow[d] \\
e \arrow[r] \arrow[ur] & X' & X \arrow[l, hook']
\end{tikzcd}

\[
  \mathrm{Hom}(e,x) \;\simeq\; \pi_0\bigl(e \backslash \Psi / x\bigr)
  \;\simeq\; \pi_0\bigl(S/x\bigr)
\]
\[
  \varinjlim_{x \in X} \pi_0\bigl(S/x\bigr)
\]

\begin{tikzcd}
S \arrow[d] \\
e \backslash X \arrow[d] \\
X
\end{tikzcd}

\note{la flèche du haut porte « cocofinal », celle du bas « cofibr. »}

\uncertain{asph.}, \uncertain{asph.}, \uncertain{cofinal}, car $X$ 0-connexe.

$S$ :

\begin{tikzcd}[column sep=small, row sep=small]
A_1 \arrow[r] \arrow[dr] & B \\
A_2 \arrow[r] & c
\end{tikzcd}

\end{document}
